\begin{sidewaystable*}[p]
\centering\footnotesize\setlength\tabcolsep{4pt}\renewcommand{\arraystretch}{1.05}
\caption{What each validity test licenses, and what it does not. No test is a prerequisite for the base reading of its family: held-out performance above the study-stated reference already licenses column~2, and each class licenses the strictly stronger statement in column~3, which also states the reading the class does \emph{not} support. A row of \S\ref{sec:findings} lacking a class is therefore \emph{narrower}, not uncontrolled, and the classes are never summed. How many records report each class, with a one-clause version of this rubric, is Table~\ref{tab:vcov} of the article. Intervention specificity applies when the operation targets a specific direction, feature, unit or token, named or not; distributional corrections, editing and unlearning are outside it and are not counted as failures, the comparison each needs being stated in Table~\ref{tab:families}.}
\label{tab:rubric}
\begin{tabular}{@{}L{3.0cm}L{3.6cm}L{10.5cm}L{3.6cm}@{}}
\toprule
Validity-test class \emph{(designs it applies to)} & What the design already supports without it & What the class adds, and what it does not & Tests coded under it (records) \\
\midrule
Sample-level permutation control \emph{(fitted probes and concept vectors; intrinsic designs)} & the fitted pipeline returns the reported value on these data & the value exceeds what the \emph{whole fitted procedure} returns when the labels are permuted across examples and it is refitted, which bounds what it can extract from labels carrying no information. It does \emph{not} hold the structure of the task fixed, which is what a randomized control task adds, nor show that the readout tracks the target rather than a correlated attribute & sample-level label permutation with refit (3) \\
Reference-distribution calibration \emph{(all decoding designs; intrinsic designs; structural and geometric)} & the statistic takes the reported value on these data & the value is placed on a scale --- an analytical chance level, a hypergeometric null, or a distribution over unrelated pairs --- so it can be read as high or low rather than as a bare number. It calibrates the statistic and does \emph{not} bound what the fitting procedure could have extracted & chance, shuffled or null reference (4) \\
Randomized control task \emph{(fitted probes and concept vectors; intrinsic designs)} & the measurement returns the reported value on these data & the readout is selective rather than expressive: a control task in the sense of Hewitt and Liang assigns each input \emph{type} a fixed random label, preserving the structure of the task while destroying its content, refits the same probe, and reports selectivity as the gap. It differs from a sample-level permutation, which reassigns labels per example and is reported in \Vperm{} records; no record reports a control task proper. The construction needs repeated input \emph{types}, which medical images seldom have and which is not coded per record, so the denominator is an \emph{upper bound on applicability} rather than a count of omitted safeguards (\S\ref{sec:future}) & none reported \\
Matched-object control \emph{(all decoding designs; intrinsic designs)} & the measurement returns the reported value on these data & the reported quantity follows the fitted object rather than any object of that size --- a direction matched in rank to the fitted one, an equal-size random unit set, a no-encoder baseline. It does \emph{not} say what the readout tracks, which is the control task above & matched-object control (1) \\
Planted ground truth \emph{(all decoding designs; intrinsic designs)} & the operator returns some value on real data, whose correctness is unknown & the operator returns the known answer on a condition built to have one (an injected bias, a planted encoder, a graded confound mixture). This validates the \emph{instrument}; it carries to the real-data claim only under the assumption that the planted condition resembles it & planted ground truth (4) \\
Learned-representation attribution \emph{(all decoding designs; intrinsic designs; geometry; tracing)} & the trained model's representation affords the reported value & the structure is a product of training rather than of the architecture and the input statistics, since an untrained or randomised encoder of the same shape does not afford it & randomly initialised encoder (4), embedding-randomisation control (1) \\
Comparison model \emph{(structural and geometric)} & the statistic takes the reported value on this model & the value differs from another model's under the same measurement --- a general encoder, a supervised classifier or a conventional feature pipeline. It does \emph{not} establish that the value is unusual in absolute terms; that needs a null & comparison model (16) \\
Semantic validation \emph{(feature decomposition and discovery)} & a unit exists in the fitted decomposition and its top activations can be displayed & the unit corresponds to the concept it is named for, as judged by clinicians or by structured labels, rather than by the analyst's reading of its exemplars & label validation (10), clinician validation (7), clinician and label validation (1) \\
Input-space corroboration \emph{(attribution, localisation and tracing)} & the map or component covaries with the output & the output depends on the region or input property the map highlights, since occluding, deleting, flipping or counterfactually altering it moves the output. It does \emph{not} establish that any particular internal component carries the behaviour: the intervention is on the input, not on the representation & perturbation confirmation (4), counterfactual confirmation (1), occlusion confirmation (1) \\
Internal-component perturbation \emph{(tracing; intrinsic designs)} & the map or component covaries with the output & the output changes when the internal component is perturbed --- masked, patched or replaced --- rather than when the input is. The credited designs are perturbations, not identified on-manifold interventions: neither on-manifold editing nor magnitude control is coded, so the class licenses causal implication under those assumptions rather than establishing them. Specificity to \emph{that} component needs a matched control & activation patching or masking (1), intervention on the concept layer (1) \\
Intervention specificity \emph{(representation intervention, targeted operations only)} & the edit changes the output & the change follows the targeted direction, feature, unit or token rather than the act of perturbing a representation of that size & matched control (3) \\
\bottomrule
\end{tabular}
\end{sidewaystable*}
